% Generated draft: requires mathematical and source review.
\documentclass[11pt]{article}
\usepackage[margin=25mm]{geometry}
\usepackage{fontspec}
\setmainfont{Latin Modern Roman}
\usepackage{amsmath,amssymb,mathtools,unicode-math}
\setmathfont{Latin Modern Math}
\usepackage{enumitem,needspace,xurl,hyperref}
\setlength{\parindent}{0pt}
\setlength{\parskip}{5pt}
\setlength{\emergencystretch}{3em}
\newcommand{\N}{\mathbb{N}}
\newcommand{\Z}{\mathbb{Z}}
\newcommand{\Q}{\mathbb{Q}}
\newcommand{\R}{\mathbb{R}}
\newcommand{\C}{\mathbb{C}}
\newcommand{\F}{\mathbb{F}}
\newcommand{\A}{\mathbb{A}}
\newcommand{\PP}{\mathbb{P}}
\newcommand{\E}{\mathbb{E}}
\newcommand{\eps}{\varepsilon}
\newcommand{\dd}{\,\mathrm{d}}
\begin{document}
\section*{UnsolvedMath ID: 701084}
\textbf{Source problem code:} catalogue entry
\subsection{Definitions and mathematical statement}
Let \(G\) be a group, and define its set of elements of finite order by
\[
T(G)=\{g\in G: g^n=1\text{ for some integer }n\geq 1\}.
\]
A subgroup \(L\leq G\) is locally finite if every subgroup generated by finitely many elements of \(L\) is finite. An element \(x\in G\) has prime order \(p\) if \(p\) is prime, \(x\neq 1\), and \(x^p=1\).

The source assumes that \(G\) is a \emph{Shunkov group}, also called a \emph{conjugately biprimitively finite group}, and refers to Kourovka Notebook Problem 6.59 for this condition. That referenced definition is absent from the supplied record. The named hypothesis is therefore retained here without reconstructing its definition; its precise scope requires review against that reference.

\textbf{Question.} For every Shunkov group \(G\) such that \(T(G)\) is infinite, is it true that
\[
\forall p\text{ prime}\;\forall x\in G,\qquad
\bigl(x\neq 1\ \text{and}\ x^p=1\bigr)
\Longrightarrow
\exists L\leq G\;
\bigl(x\in L,\ L\text{ is infinite and locally finite}\bigr)?
\]
The hypothesis requires infinitely many elements of finite order; it does not require every element of \(G\) to have finite order.

\textbf{Supplied qualification.} The source reports that the conclusion holds for a fixed element \(x\) of prime order under the additional hypothesis
\[
\forall g,h\in G,\qquad
\langle g^{-1}xg,h^{-1}xh\rangle\text{ is soluble}.
\]
Here a group is soluble if its derived series reaches the trivial subgroup after finitely many steps. The source attributes this additional-case result to V. P. Shunkov's book cited below; it is recorded as a supplied assertion, not an independently verified result.
\subsection{Short English statement}
The problem asks whether every element of prime order in a Shunkov group with infinitely many elements of finite order lies in an infinite locally finite subgroup. The source reports an affirmative answer for an element whose every pair of conjugates generates a soluble subgroup. The definition of a Shunkov group is deferred to an unavailable cross-reference.
\subsection{Sources}
\begin{itemize}
\item Local UnsolvedMath registry, supplied record with permanent ID \(701084\), exact source code \texttt{catalogue entry}, title \emph{Kourovka Notebook Problem 14.100}; upstream revision \texttt{d371ee917ac9647fbb06cd6fc529deec6f3dda82}. The supplied registry status is \texttt{open}; this is a source label, not an independently verified literature assessment.
\item \emph{The Kourovka Notebook (Unsolved Problems in Group Theory)}, Problem 14.100, proposed by A. K. Shl\"epkin in the 14th issue (1999). The record also identifies the 21st issue edition (2026), \texttt{arXiv:1401.0300v48}. The supplied collection description names E. I. Khukhro and V. D. Mazurov as editors. Problem 6.59 is the supplied cross-reference for the Shunkov-group hypothesis; its text is not included.
\item V. P. Shunkov, \emph{\(M_p\)-groups}, Moscow, Nauka, 1990 (Russian), cited by the supplied statement for the additional case involving soluble subgroups generated by pairs of conjugates.
\end{itemize}
\end{document}
